\documentclass[10pt]{article}
% Math Packages
\usepackage{amsmath, mathtools}
\usepackage{amssymb}
\usepackage{amsthm}
\usepackage{amsfonts}
\usepackage{bbm}
\usepackage{breqn}
\usepackage[margin=1in]{geometry}
\usepackage{graphicx}
\usepackage{tikz}
\usetikzlibrary{arrows.meta}
\usetikzlibrary{calc}
\usepackage{forest}
\usepackage{tikz-qtree}
\graphicspath{ {./images/} }
\usepackage{hyperref}
\usepackage[capitalize]{cleveref}
\usepackage[shortlabels]{enumitem}
\usetikzlibrary{arrows,matrix,positioning}
\usepackage{multicol}

% for the pipe symbol
\usepackage[T1]{fontenc}

% Citing theorems by name. (source: https://tex.stackexchange.com/questions/109843/cleveref-and-named-theorems)
\makeatletter
\newcommand{\ncref}[1]{\cref{#1}\mynameref{#1}{\csname r@#1\endcsname}}

\def\mynameref#1#2{%
  \begingroup
  \edef\@mytxt{#2}%
  \edef\@mytst{\expandafter\@thirdoffive\@mytxt}%
  \ifx\@mytst\empty\else
  \space(\nameref{#1})\fi
  \endgroup
}
\makeatother

% Colorful Notes
\usepackage{color} \definecolor{Red}{rgb}{1,0,0} \definecolor{Blue}{rgb}{0,0,1}
\definecolor{Purple}{rgb}{.5,0,.5} \def\red{\color{Red}} \def\blue{\color{Blue}}
\def\gray{\color{gray}} \def\purple{\color{Purple}}
\newcommand{\rnote}[1]{{\red [#1]}} % \rnote{foo} gives '[foo]' in red
\newcommand{\pnote}[1]{{\purple [#1]}} % \pnote{foo} gives '[foo]' in purple
\newcommand{\bnote}[1]{{\blue #1}} % \bnote{foo} gives 'foo' in blue
\newcommand{\gnote}[1]{{\gray #1}} % \gnote{foo} gives 'foo' in gray
\newcommand{\Max}[1]{{\purple [#1]}} % \bnote{foo} then 'foo' is blue


% Claim numbering (the counter restarts after each proof environment)
\newcounter{claimcount}
\setcounter{claimcount}{0}
\newenvironment{claim}{\refstepcounter{claimcount}\par\addvspace{\medskipamount}\noindent\textbf{Claim \arabic{claimcount}:}}{}
\usepackage{etoolbox}
\AtBeginEnvironment{proof}{\setcounter{claimcount}{0}}
\newenvironment{claimproof}{\par\addvspace{\medskipamount}\noindent\textit{Proof of Claim  \arabic{claimcount}.}}{\hfill\ensuremath{\qedsymbol} \tiny{Claim}

  \medskip}
% Add claim support to cleverref
\crefname{claimcount}{Claim}{Claims}


% Math Environments
\newtheorem{theorem}{Theorem}
\newtheorem{assumption}[theorem]{Assumption}
\newtheorem{lemma}[theorem]{Lemma}
\newtheorem{proposition}[theorem]{Proposition}
\newtheorem{corollary}[theorem]{Corollary}
\newtheorem{question}[theorem]{Question}
\theoremstyle{definition}
\newtheorem{definition}[theorem]{Definition}
\newtheorem{remark}[theorem]{Remark}
\newtheorem{example}[theorem]{Example}
\newtheorem{notation}[theorem]{Notation}
\newtheorem{problem}[theorem]{Problem}

% Matrices and Column Vectors. 
\usepackage{stackengine}
\setstackgap{L}{1.0\normalbaselineskip}
\usepackage{tabstackengine}
\setstackEOL{;}% row separator
\setstackTAB{,}% column separator
\setstacktabbedgap{1ex}% inter-column gap
\setstackgap{L}{1.5\normalbaselineskip}% inter-row baselineskip
\let\nmatrix\bracketMatrixstack  %Usage: \nmatrix{1,2,3\4,5,6}
\newcommand\cv[1]{\setstackEOL{,}\bracketMatrixstack{#1}} %usage: \cv{1,2,3}

% Custom Math Coqmmands
\newcommand{\vt}{\vskip 5mm} % vertical space
\newcommand{\fl}{\noindent\textbf} % first line
\newcommand{\Fl}{\vt\noindent\textbf} % first line with space above
\newcommand{\norm}[1]{\left\lVert#1\right\rVert} % norm
\newcommand{\pnorm}[1]{\left\lVert#1\right\rVert_p} % p-norm
\newcommand{\qnorm}[1]{\left\lVert#1\right\rVert_q} % q-norm
\newcommand{\1}[1]{\textbf{1}_{\left[#1\right]}} % indicator function 
\def\limn{\lim_{n\to\infty}} % shortcut for lim as n-> infinity
\def\sumn{\sum_{n=1}^{\infty}} % shortcut for sum from n=1 to infinity
\def\sumkn{\sum_{k=1}^{n}} % shortcut for sum from k=1 to n
\def\sumin{\sum_{i=1}^{n}} % shortcut for sum from i=1 to n
\def\SAs{\sigma\text{-algebras}} % shortcut for $\sigma$-algebras
\def\SA{\sigma\text{-algebra}} % shortcut for $\sigma$-algebra
\def\Ft{\mathcal{F}_t} % time-indexed sigma-algebra (t)
\def\Fs{\mathcal{F}_s} % time-indexed sigma-algebra (s)
\def\F{\mathcal{F}} % sigma-algebra
\def\G{\mathcal{G}} % sigma-algebra
\def\R{\mathbb{R}} % Real numbers
\def\N{\mathbb{N}} % Natural numbers
\def\Z{\mathbb{Z}} % Integers
\def\E{\mathbb{E}} % Expectation
\def\P{\mathbb{P}} % Probability
\def\Q{\mathbb{Q}} % Q probability
\def\dist{\text{dist}} %Text 'dist' for things like 'dist(x,y)'
\newcommand{\indep}{\perp \!\!\! \perp}  %independence symbol
\def\Var{\mathrm{Var}} % Variance
\def\tr{\mathrm{tr}} % trace

% Brackets and Parentheses
\def\[{\left [}
    \def\]{\right ]}
% \def\({\left (}
%   \def\){\right )}



\usepackage{color}
\definecolor{Red}{rgb}{1,0,0}
\definecolor{Blue}{rgb}{0,0,1}
\definecolor{Purple}{rgb}{.5,0,.5}
\def\red{\color{Red}}
\def\blue{\color{Blue}}
\def\gray{\color{gray}}
\def\purple{\color{Purple}}
\definecolor{RoyalBlue}{cmyk}{1, 0.50, 0, 0}
\newcommand{\dempfcolor}[1]{{\color{RoyalBlue}#1}} 
\DeclareRobustCommand{\demph}[1]{% definitions
  \ifmmode
  {\color{Blue}#1}%
  \else
  {\color{Blue}\bfseries\slshape #1}%
  \fi
}
% comment exactly one of the following line to show / hide the solutions
% \newcommand{\solution}[1]{{\purple #1}} % uncomment to show the solutions
\newcommand{\solution}[1]{} % uncomment to hide the solutions



\title{Lecture Notes for Math 372: \\Elementary Probability and Statistics}
\date{Last updated: \today}
% \author{mh}

\begin{document}
\begin{center}
  \section*{Math 372: Homework 06}
  \textit{Due Monday, October 12}
\end{center}


\begin{problem}[Variance shortcut formula]
  Let $X$ be a random variable, and let $\mu= \E\left[X \right]$. By
  definition of variance, we know that
  $\Var(X) = \E\left[(X-\mu)^{2} \right].$ Show that $\Var(X) = \E\left[X^{2} \right] - \mu^{2}.$
\end{problem}


\begin{problem}  
  \label{ex:prussian-officers-killed-horsekick}
  Let $K$ be the number of Prussian cavalry officers killed by horsekick in a random year. Because of
  careful data collection by the Prussian military\footnote{This problem actually comes from real
    data collected by the Prussian army, and was one of the first known applications of the Poisson
    distribution, by Russian mathematician L. Bortkewitsch in 1898; see
    \url{https://www.jstor.org/stable/2348169?seq=3}}, we know that $K$ has the following distribution:
  \begin{equation*}
    p(0) = .51, \quad p(1) = .33, \quad p(2) = .11, \quad p(3) = .04, \quad p(4) = .01, \quad \text{and} \quad  \P\left[K \geq 5\right] =0.
  \end{equation*}
  
  \begin{enumerate}[(a)]
    \item What is the expectation of the random variable $K$? %(Ans: $\approx 0.71$). 
    \item What is the variance of $K$? %(Ans: $\approx 0.79$)
    \item Let $\lambda= \E\left[K \right]$, rounded to 1 decimal place, which
    you computed in part $(a)$. Using this value, let $Y$ be a Poisson random
    variable with rate parameter $\lambda$, and let
    $q(y) := \P\left[Y=y \right]$ be the pmf of $Y$. Compute
    $q(0),q(1),q(2),q(3),q(4),$ and $\P\left[Y \geq 5\right]$, and compare the
    probability mass functions $p$ and $q$.
  \end{enumerate}
\end{problem}


\begin{problem}
  Given an example with dice that show that, in general,
  \begin{equation*}
    \P(B) \neq \P(B|A) + \P(B|A^{c}).
  \end{equation*}
\end{problem}



\begin{problem}
  Compute $\E\left[X \right]$, $\E\left[X^{2} \right]$ and $\Var(X)$ for the
  following random variables:
  \begin{enumerate}[(a)]
    \item $X = (\text{value of one roll of a 6-sided dice})$.
    \item A fair coin toss.
    $Y = \left\{ \begin{array}{l@{\quad:\quad}l} 1 & \text{with probability
        }\frac{1}{2}\\0 & \text{with probability
        }\frac{1}{2} \end{array}\right.$
    \item A biased coin toss. $Z = \left\{ \begin{array}{l@{\quad:\quad}l} 1 &  \text{with
          probability }p\\0 & \text{with probability }1-p \end{array}\right.$
  \end{enumerate}
\end{problem}


\begin{problem}%[independence]
  Alice rolls six dice and wins if she rolls at least one six. Bob rolls
  twelve dice and wins if he rolls at least two sixes. Who has the greater
  probability of winning? \textit{[Hint: calculate the probabilities to lose.]}
\end{problem}



\begin{problem}
  A certain brand of upright freezer is available in three different rated
  capacities: $16$,  $18$, and $20$ cubic feet. Let $X$ be the rated capacity
  of a freezer of this brand sold at a certain store. Suppose that $X$ has
  probability mass function
  \begin{center}
    \begin{tabular}[c]{c|ccc}
      $x$&16&18&20\\ 
      \hline
      $p(x)$&.2&.5&.3 
    \end{tabular}
  \end{center}
  \begin{enumerate}[(a)]
    \item Compute $\E\left[X \right], \E\left[X^{2} \right], \Var(X)$.
    \item If the price of a freezer having capacity $X$ is $70X-650$, what is
    the expected price paid by the next customer to buy a freezer?
    \item What is the variance of the price paid by the next customer?
    \item Suppose that although the rated capacity of a freezer is $X$, the
    actual capacity is $h(X)=X-.008X^{2}$. What is the expected actual
    capacity of the freezer purchased by the next customer.
  \end{enumerate}
\end{problem}

\begin{problem}[Exploding Dice]
  Compute the expectation of an exploding 6-sided dice.
  \textit{\textbf{Note: }An `exploding' dice roll is a mechanic used in various tabletop games whereby the player rolls a dice,
  and then if they rolled the maximum number on the die, they get to roll it again and add the new roll to
  the old one. This can happen more than once.} 
\end{problem}




\begin{problem}[Geometric random variable]
  A basketball player is shooting hoops. Suppose that for each shot, the
  probability of successfully making the shot is $p\in (0,1)$, (where $p$
  depends on the skill of the player), and the probability of missing the shot
  is $1-p$. Let
  \begin{equation*}
    X = \text{(\# of throws until they make a basket)}
  \end{equation*}
 
  \begin{enumerate}[(a)]
    \item Explain why the pmf of $X$ is
    \begin{equation*}
      \P\left[X=k \right] = \left\{ \begin{array}{l@{\quad:\quad}l} (1-p)^{k-1}p & k=1,2,\ldots\\0 &\text{otherwise} \end{array}\right.
    \end{equation*}
    (A random variable with this pmf is called a \demph{geometric random
      variable} with success probability $p$, and we write
    $X\sim \operatorname{Geom}(p)$. The parameter $p$ is called the
    \demph{success probability}.)
    \item Compute $\P\left[X \leq k\right]$ for $k=1,2,\ldots$. \textit{Hint:} first
    compute $\P\left[X>t \right]$.
    \item Plot the cdf $F$ of $X$ for $p=1/2$. (Note: it should be a
    nondecreasing step function with $F(-\infty)=0$ and $F(+\infty)=1$.)
    \item Suppose that $X$ is a geometric random
    variable with success parameter $p\in (0,1)$. Show that 
    \begin{equation*}
      \P\left[X=m+k \mid X>m\right] = \P\left[X=k \right] \quad \text{for
        integers $m>0$ and $k>1$.}
    \end{equation*}
    Explain why this is called the \textit{memoryless property} of the geometric distribution.
    % \item Now suppose that $T\sim \exp(\lambda)$ (in other words, $T$ is an exponential random variable
    % rate $\lambda>0$). Show that
    % \begin{equation*}
    %   \P\left[ T>s+t\mid T>s\right] = \P\left[T>t \right]
    % \end{equation*}
    % and interpret this result.
  \end{enumerate}
\end{problem}


\begin{problem}
  The human genome consists of 3.2 billion nucleotides
  $(\mathtt{A},\mathtt{G},\mathtt{C},\mathtt{T})$. For this problem we can
  think of the human genome as a sequence of nucleotides
  \begin{center}
    \texttt{CAGGAAGCATTTATATAGAGGTAATAATTAAATTTAACTTTTTATTATTCTCTATGCCCAAGAGAATGTG...}
  \end{center}
  consisting of 3,200,000,000 nucleotide letters. The DNA mutation rate in
  humans is approximately $2\times 10^{-8}=0.00000002$ mutations per
  nucleotide per generation. Assume that nucleotides mutate
  \textit{independently}, that is, whether one nucleotide undergoes mutates is
  independent from any other nucleotides.

  
  Let $X$ be the number of mutations in
  the genome of a randomly selected newborn baby.
  \begin{enumerate}[(a)]
    \item Explain why $X$ is a binomial random variable.
    Identify the value of the two parameters $n$
    (number of trials) and $p$ (success probability).
    \item \label{item:3} The average number of mutations in a newborn baby's genome is
    $\mu:=\E\left[X \right]$. Find $\mu$.
    \item  For statistical
    reasons that we will get to later in the class,\footnote{The Central Limit
      Theorem} over 95\% of babies will have a number of mutations in the
    interval
    \begin{equation*}
      (\mu-\sigma,\mu+\sigma)
    \end{equation*}
    where $\sigma = \sqrt{\Var(X)}$. Find $\sigma$ and (using your answer to
    \ref{item:3}) find the interval $(\mu-\sigma,\mu+\sigma)$.
    \item Attempt to compute some probabilities, such as $\P\left[X=1 \right]$ and $\P\left[X=60
    \right]$, and $\P\left[X \leq 100\right]$. I do not expect you to be successful with this. What
    makes this task difficult?
    \item Let $Y$ be a Poisson
    distribution with mean $\mu$. Explain why $Y$ can be used to approximate
    $X$, and compute the following probabilities:
    \begin{enumerate}[label=(\roman*)] 
      \item $\P\left[Y=0 \right]$
      \item $\P\left[Y=1 \right]$
      \item $\P\left[Y=5 \right]$
      \item $\P\left[Y \geq 3\right]$
    \end{enumerate}
  \end{enumerate}
\end{problem}



\begin{problem}
  Let $Y$ be a Poisson random variable with parameter $\lambda$. Find $\E\left[\frac{1}{Y+1} \right]$.
\end{problem}



% \begin{problem}[Properties of the geometric distribution]
%   Let $\alpha\in (0,1)$, and let $X$ be a random variable with probability mass function
%   \begin{equation*}
%     p(x) = \left\{ \begin{array}{l@{\quad}l} \left( 1-\alpha \right)^{x-1}\alpha  & \text{if }x=1,2,3\ldots\\0 & \text{otherwise.} \end{array}\right.
%   \end{equation*}
%   (A random variable with this probability mass function is called a \textit{geometric random variable} with success parameter $\alpha$.)

%   \begin{enumerate}[(a)]
%     \item We know that in general probability mass functions should always sum
%     up to one. Let's check that it is indeed the case here. Specifically,
%     carefully show that $p(1)+p(2)+p(3)+\ldots =1$, making sure to justify
%     your steps.
%     \item Next we will compute the expected value of $X$ using the formula 
%     \begin{equation*}
%       \E\left[X \right] := \sum_{x}x p(x),
%     \end{equation*}
%     where $x$ ranges over all possible values of $X$.
%     \begin{enumerate}[(i)]
%       \item \label{item:1} Expand out above summation, including the first five terms followed by ``+\dots''
%       \item \label{item:2} The geometric series formula tells us that as long as $|x|<1$, the following
%       formula holds:
%       \begin{equation*}
%         a + ax+ax^{2}+ax^{3}+ax^{4}+ax^{5}+\ldots = \frac{a}{1-x}.
%       \end{equation*}
%       Differentiate both sides of the above equation with respect to $x$. 
%       \item Use your answers to parts \ref{item:1} and\ref{item:2} to deduce the value of
%       $\E\left[X \right]$
%       \item What is the value of $\E\left[X \right]$ if $\alpha=\frac{1}{2}$? What about if
%       $\alpha=\frac{2}{3}$?
%     \end{enumerate}
%   \end{enumerate}
% \end{problem}


% \begin{problem}
%   Consider another approach where we show that
%   \begin{equation*}
%     \E\left[e^{tX} \right] = \exp \left( \lambda \left( e^{t}-1 \right)  \right) 
%   \end{equation*}
%   and then differentiate once or twice, and then set $t=0$. This would be a nice problem to introduce mgfs
% \end{problem}








% \begin{problem}[Extra credit]
%   Suppose you draw a pair of cards from a standard deck of cards. Let $B$ denote the event that both
%   cards are aces. Let $A$ be the event that you have at least one ace. Let $S$ be the event that you
%   have the ace of spaces. In Homework 4, we saw that
%   \begin{equation*}
%     \P\left[B\mid A \right] = \frac{1}{33} \approx 0.03
%   \end{equation*}
%   and that
%   \begin{equation*}
%     \P\left[B\mid S \right] = \frac{1}{17}\approx 0.06
%   \end{equation*}
%   This is such a strange result! Why is the second one twice as big as the first? Could we have made a
%   mistake?

%   Write program using any language that estimates the conditional probability $\P\left[B\mid A \right]$.
%   Also write a program that estimates conditional probability $\P\left[B\mid S \right]$. What
%   estimates do you get?
% \end{problem}


\end{document}

